% The LINEAR (least-squares nestedLSQ) semi-Lagrangian level-set method as an
% INDEPENDENT method line: its own studies (config/linearConv*.yaml), its own
% reveal deck (../lsl-level-set-presentation/lsl.template.html) and this
% article. Single data source: data/{figures,tables} is written by the
% snakemake pipeline (make sl-linear) and mirrored verbatim into the
% presentation's data/ (dual-copy, workflow/scripts/propagate_data.py), so the
% deck and this paper always show the same, newest results.
%
% Figures/tables are guarded with \IfFileExists so the article compiles before
% the first study run and simply grows as results land.
\documentclass[preprint,3p,11pt]{elsarticle}

\usepackage[T1]{fontenc}
\usepackage{amsmath,amssymb}
\usepackage{graphicx}
\usepackage{booktabs}
\usepackage{url}
\usepackage[hidelinks]{hyperref}

\graphicspath{{data/figures/}}

% notation shared with the companion quadratic semi-Lagrangian article
\newcommand{\psilev}{\psi}
\newcommand{\uvec}{\mathbf{u}}
\newcommand{\xc}{\mathbf{x}_c}
\newcommand{\xd}{\mathbf{x}_d}
\newcommand{\gc}{\mathbf{g}_c}

% guarded figure/table inputs: compile cleanly before the data exists
% (\path renders the raw file names -- underscores included -- in the fallback)
\newcommand{\guardedgraphics}[2][width=\linewidth]{%
  \IfFileExists{data/figures/#2}{\includegraphics[#1]{#2}}{%
    \fbox{\parbox{0.9\linewidth}{\centering\small pending pipeline output:\\
      \path{data/figures/#2}\\(run \texttt{make sl-linear})}}}}
\newcommand{\guardedinput}[1]{%
  \IfFileExists{#1}{\input{#1}}{%
    \fbox{\parbox{0.9\linewidth}{\centering\small pending pipeline output:\\
      \path{#1}\\(run \texttt{make sl-linear})}}}}

\begin{document}
\begin{frontmatter}

\title{A linear semi-Lagrangian level-set method on unstructured meshes:
least-squares reconstruction, and why the raw Taylor extrapolation is not
enough}

\author[tuda]{Tomislav Mari\'{c}}
\address[tuda]{Mathematical Modeling and Analysis, TU Darmstadt, Germany}

\begin{abstract}
We document the linear semi-Lagrangian level-set method as an independent
method line next to its higher-order companion, the uncached quadratic
weighted least-squares reconstruction (UQWLSR). All three reconstructions of
the framework share the same backward, second-order Taylor departure foot and
the same tetrahedral-decomposition phase indicator; they differ in exactly
one ingredient --- how the old-time level set is reconstructed at the
departure foot. We study the two \emph{linear} choices and find them to be
qualitatively opposite. The raw first-order Taylor extrapolation
(\texttt{linearTaylor}), which evaluates the arrival cell's own gradient at
the foot, is unstable in pure advection: the signed-distance gradient defect
$\||\nabla\psi|-1\|$ grows geometrically under mesh refinement (to
$\mathcal{O}(10^{9})$ at $N=128$ on the reversed vortex) and even the
interface shape error mildly diverges, because the single-cell extrapolation
has no mechanism to damp the steepening it induces. The least-squares linear
reconstruction (\texttt{nestedLSQ}), which fits a linear polynomial over the
neighbour stencil, is stable and convergent: the stencil averaging removes
the runaway, and the shape, volume and \emph{band} gradient errors all
decrease with refinement. We therefore build the linear method line on
\texttt{nestedLSQ} and verify it on the reversed 2D vortex and the 3D shear
and deformation benchmarks, hexahedral and polyhedral, reporting the
convergence of \emph{all} error measures together with volume-fraction field
visualizations. The method is convergent but only \emph{conditionally}
stable, and its safe resolution is flow-dependent: it converges to $256^2$
(2D vortex) and to $N=160$ (3D deformation, both mesh types) at
$\mathrm{CFL}=\tfrac12$, but the strongest sustained 3D shear destabilizes it
on hexahedra past $N=80$ --- a limit the polyhedral face-neighbour stencil
pushes beyond reach. The unstable \texttt{linearTaylor} variant is documented
as a characterized limitation --- and as the reconstruction that nonetheless
serves the two-phase solver, where its dissipation and the short integration
times keep it in its stable regime. The picture is two design points of one
framework plus one cautionary tale: \texttt{nestedLSQ} the robust, low-cost
linear transport method; UQWLSR the kinematic-accuracy champion;
\texttt{linearTaylor} a reconstruction that must never be used for
reinitialization-free advection.
\end{abstract}

\begin{keyword}
level set \sep semi-Lagrangian advection \sep unstructured meshes
\sep OpenFOAM \sep least-squares reconstruction
\end{keyword}
\end{frontmatter}

% ============================================================================
\section{Introduction}
\label{sec:intro}

The level-set method \cite{osher1988,sethian1999,sussman1994} transports a
scalar field $\psilev$ whose zero set is the fluid interface. On unstructured
finite-volume meshes, Eulerian discretizations of the level-set equation
couple the interface accuracy to flux limiters and to the divergence scheme;
semi-Lagrangian schemes \cite{staniforth1991,strain1999} avoid this
altogether by tracing characteristics backward and reconstructing the
old-time field at the departure foot. In our companion work, a per-cell
quadratic weighted least-squares reconstruction (UQWLSR) provides a
high-order reconstruction on arbitrary polyhedral stencils; that line
delivers the best kinematic accuracy of the framework.

This article documents the \emph{linear} reconstruction --- the natural,
low-cost sibling of the quadratic fit --- as an independent method line, and
in doing so reports a result that is easy to get wrong. There are two ways to
reconstruct a linear approximation of $\psilev$ at the departure foot, and on
a reinitialization-free semi-Lagrangian scheme they behave in opposite ways.
The obvious one, which we call \texttt{linearTaylor}, evaluates the arrival
cell's own gradient as a first-order Taylor expansion at the foot; it needs a
single cell-gradient evaluation and nothing else. We find it to be
\emph{unstable} in pure advection: the signed-distance gradient defect
$\bigl\||\nabla\psilev|-1\bigr\|$ grows geometrically under mesh refinement
(to $\mathcal{O}(10^{9})$ at $N=128$ and $\mathcal{O}(10^{21})$ at $N=256$ on
the reversed vortex) and even the interface shape error mildly diverges,
because a single-cell extrapolation across the foot displacement has no
mechanism to damp the gradient steepening it induces, and the reversed-vortex
shear steepens gradients relentlessly. The stencil-bounds clamp does not cure
it. The second choice, \texttt{nestedLSQ}, fits a linear polynomial over the
neighbour stencil by least squares --- the same stencil averaging that makes
the quadratic fit well behaved --- and is stable and convergent: shape,
volume and the band-restricted gradient defect all decrease with refinement.

We therefore build the linear method line on \texttt{nestedLSQ}. It keeps the
attractions that motivate a linear scheme --- a compact least-squares solve
with no cached per-cell pseudo-inverse, $O(1)$ memory per cell, and no
fine-mesh memory ceiling --- while remaining a genuine, convergent transport
method. The unstable \texttt{linearTaylor} variant is not discarded but
documented (Section~\ref{sec:lintaylor}): it is the reconstruction used in the
two-phase ``consistent-linear'' pipeline, where the coupled dissipation and
the short integration times before capillary equilibrium keep it inside its
stable regime, and where it was measured to outlive the quadratic pipeline at
every resolution. Its pure-advection instability and its two-phase utility are
both real and not in tension: the coupled comparison is relative (linear less
destabilizing than quadratic in the curvature-feedback loop), the kinematic
result absolute (unbounded gradient growth over a long reversible integration).

The contributions here are: (i) a complete description of the linear method
--- departure foot, least-squares linear reconstruction, safeguards, phase
indicator --- and a characterization of why the raw Taylor extrapolation
fails; (ii) verification of \texttt{nestedLSQ} on the standard 2D and 3D
kinematic benchmarks \cite{rider1998,leveque1996,enright2002,liovic2006} on
hexahedral \emph{and} polyhedral \cite{juretic_cfmesh} meshes with the
convergence of \emph{all} error measures, not only the headline shape error;
and (iii) volume-fraction field visualizations that make the method's
behaviour and its resolution dependence directly visible.

% ============================================================================
\section{Method}
\label{sec:method}

\subsection{Level-set transport and semi-Lagrangian update}
The interface $\Sigma(t)$ is the zero set of $\psilev(\mathbf{x},t)$,
transported by
\begin{equation}
  \partial_t \psilev + \uvec\cdot\nabla\psilev = 0 ,
  \label{eq:levelset}
\end{equation}
which keeps $\psilev$ constant along characteristics
$\dot{\mathbf{x}}=\uvec$. Tracing the characteristic that arrives at a cell
centre $\xc$ at $t^{n+1}$ backward over one time step to its departure foot
$\xd$ turns conservation-along-trajectories into an exact one-step update,
\begin{equation}
  \psilev^{\,n+1}(\xc) = \psilev^{\,n}(\xd),
  \label{eq:sl}
\end{equation}
evaluated at every cell centre. No flux, no divergence scheme and no linear
system appear; the scheme has exactly two ingredients --- locating $\xd$
(Section~\ref{sec:foot}) and reconstructing $\psilev^{\,n}$ there
(Section~\ref{sec:recon}).

\subsection{Second-order departure foot}
\label{sec:foot}
The foot is located by the backward Taylor expansion of the arriving
characteristic,
\begin{equation}
  \xd = \xc - \uvec^{\,n+1}\,\Delta t
        + \tfrac12\!\left[\frac{\uvec^{\,n+1}-\uvec^{\,n}}{\Delta t}
                 + \left(\uvec^{\,n+1}\!\cdot\nabla\right)\uvec^{\,n+1}\right]\!\Delta t^{2},
  \label{eq:foot}
\end{equation}
in which the bracket approximates the material acceleration
$\mathrm{D}\uvec/\mathrm{D}t$ with a one-sided temporal difference and a
Gauss-linear cell gradient. The dropped terms are $O(\Delta t^{3})$ per
step, so the foot is second-order accurate in time; this foot is shared
verbatim with the quadratic companion line and is held fixed throughout. A
foot-radius guard warns when $|\xd-\xc|$ exceeds the neighbour-stencil
radius, i.e.\ when the operational Courant limit ($\mathrm{CFL}\le 1$) is
violated.

\subsection{Least-squares linear reconstruction (\texttt{nestedLSQ})}
\label{sec:recon}
The linear method reconstructs the old-time field at the foot from a linear
polynomial fitted over the arrival cell's neighbour stencil $\mathcal{S}_c$ by
weighted least squares --- the linear-basis restriction of the quadratic fit
of the companion line:
\begin{equation}
  R_c(\mathbf{x}) = \psilev^{\,n}_c + \mathbf{g}_c\cdot(\mathbf{x}-\xc),
  \qquad
  \mathbf{g}_c = \arg\min_{\mathbf{g}} \sum_{i\in\mathcal{S}_c}
      w_i^{2}\bigl(\psilev^{\,n}_c + \mathbf{g}\cdot\mathbf{d}_i - \psilev^{\,n}_i\bigr)^2 ,
  \label{eq:nestedlsq}
\end{equation}
with $\mathbf{d}_i=\mathbf{x}_i-\xc$ and inverse-distance weights
$w_i=1/\lvert\mathbf{d}_i\rvert$; the constant term is dropped so
$R_c(\xc)=\psilev^{\,n}_c$ exactly, and $\psilev^{\,n}(\xd)\approx R_c(\xd)$.
The gradient $\mathbf{g}_c$ is thus the \emph{least-squares} slope over the
whole stencil rather than a single-cell difference; that stencil averaging is
the decisive difference from the raw Taylor extrapolation of
Section~\ref{sec:lintaylor}. Three properties follow.

\paragraph{Accuracy.} The fit is exact for linear fields, so the per-step
reconstruction error is $O(h^{2})$ (the quadratic remainder is dropped). Over
$T/\Delta t \propto 1/h$ steps the accumulated transport error is first order
in $h$; the measured 2D shape order is $\approx1.1$ and the volume order
$\approx1.5$ (Section~\ref{sec:conv2d}), well below the quadratic line's
$\approx3$ but a genuine, monotone convergence.

\paragraph{Stability and its CFL limit.} Because the slope is a stencil
average, the reconstruction does not amplify the gradient steepening that the
reversed-vortex shear induces, and the band gradient defect
$\bigl\||\nabla\psilev|-1\bigr\|$ \emph{decreases} with refinement at
$\mathrm{CFL}=\tfrac12$. The method is, however, only \emph{conditionally}
stable: at $\mathrm{CFL}=1$ it converges on the coarse meshes but
destabilizes on the finest (Section~\ref{sec:conv2d}), unlike the quadratic
line, whose order is CFL-independent up to the stability bound. All
production linear studies therefore run at $\mathrm{CFL}\le\tfrac12$.

\paragraph{Cost and memory.} The per-cell solve is a single small symmetric
system ($2\times2$ in 2D, $3\times3$ in 3D) assembled and factorised on the
fly; there is no cached pseudo-inverse, so the memory footprint is $O(1)$ per
cell and the fine-mesh memory ceiling of the cached quadratic reconstruction
does not exist. The reconstruction is local: it uses only the arrival cell's
stencil, so the only parallel communication is the standard halo exchange of
$\psilev^{\,n}$.

\subsection{The \texttt{linearTaylor} variant and its instability}
\label{sec:lintaylor}
The obvious alternative linear reconstruction is the arrival cell's own
first-order Taylor expansion,
\begin{equation}
  R_c^{\mathrm{T}}(\xd) = \psilev^{\,n}_c + (\nabla\psilev^{\,n})_c\cdot(\xd-\xc),
  \label{eq:lintaylor}
\end{equation}
with $(\nabla\psilev^{\,n})_c$ a point-cell least-squares \emph{cell} gradient
(\texttt{gradPsi pointCellsLeastSquares}). It looks cheaper still --- one
gradient field, one dot product --- but on the reinitialization-free scheme it
is \emph{unstable in pure advection}. Table~\ref{tab:lintaylor} contrasts the
two linear reconstructions on the reversed 2D vortex at $\mathrm{CFL}=\tfrac12$:
where \texttt{nestedLSQ} holds the band gradient defect bounded and decreasing,
\texttt{linearTaylor} lets it grow geometrically, reaching $\mathcal{O}(10^{9})$
at $N=128$ and $\mathcal{O}(10^{21})$ at $N=256$, and the shape error grows too
(fitted order $-0.9$). The mechanism is a positive feedback: the single-cell
extrapolation across the foot displacement overshoots where $|\nabla\psilev|$
is already too steep, steepening it further, with no stencil averaging to damp
it; finer meshes take more steps, so the runaway is worse, not better. The
stencil-bounds clamp (Section~\ref{sec:safeguards}) does not cure it --- the
reconstructed \emph{value} near the interface stays within the stencil range
while the \emph{gradient} explodes in the field around it.

\begin{table}[t]
  \centering
  \caption{The two linear reconstructions on the reversed 2D vortex,
  $\mathrm{CFL}=\tfrac12$: band gradient defect
  $\bigl\||\nabla\psilev|-1\bigr\|^{\mathrm{band}}$ and interface shape error at
  $t=T$. \texttt{nestedLSQ} converges; \texttt{linearTaylor} diverges
  geometrically. Measured; reproduce by setting \texttt{reconstruction
  linearTaylor} in \texttt{config/linearConv2Dvortex.yaml}.}
  \label{tab:lintaylor}
  \begin{tabular}{lcccc}
  \toprule
   & \multicolumn{2}{c}{band gradient defect} & \multicolumn{2}{c}{shape error} \\
  $N$ & \texttt{nestedLSQ} & \texttt{linearTaylor} & \texttt{nestedLSQ} & \texttt{linearTaylor} \\
  \midrule
  32  & $3.6\times10^{-2}$ & $1.1\times10^{1}$  & $9.2\times10^{-4}$ & $5.4\times10^{-3}$ \\
  64  & $5.8\times10^{-3}$ & $2.9\times10^{3}$  & $4.2\times10^{-4}$ & $9.9\times10^{-3}$ \\
  128 & $1.5\times10^{-3}$ & $1.4\times10^{9}$  & $1.6\times10^{-4}$ & $2.2\times10^{-2}$ \\
  256 & $4.8\times10^{-4}$ & $2.0\times10^{21}$ & $9.2\times10^{-5}$ & $3.1\times10^{-2}$ \\
  \bottomrule
  \end{tabular}
\end{table}

The variant is not merely a cautionary tale: \texttt{linearTaylor} is the
reconstruction of the two-phase ``consistent-linear'' pipeline, where it was
measured to outlive the quadratic pipeline at every resolution and to break
its finer-blows-first blow-up trend. There is no contradiction with the
instability reported here. That coupled result is \emph{relative} --- linear
less destabilizing than quadratic in the curvature-feedback loop, over the
short integration before capillary equilibrium --- while the kinematic result
is \emph{absolute}, an unbounded gradient growth over a long reversible
integration. Under two-phase coupling the solver's own dissipation and the
short time horizon keep \texttt{linearTaylor} inside its stable regime; in pure
transport to full vortex reversal, nothing does. For pure advection the linear
method of record is therefore \texttt{nestedLSQ}.

\subsection{Safeguards}
\label{sec:safeguards}
Two optional safeguards live in the reconstruction base class and are shared
by every reconstruction of the framework. (i) A Barth--Jespersen-type slope
limiter on the reconstruction increment; measured on the benchmarks it costs
accuracy without improving robustness for the linear method, and it is OFF
by default. (ii) \texttt{clipToStencilBounds}, a monotone clamp of the
reconstructed value to the stencil's $[\min,\max]$; it is required on
polyhedral meshes --- where irregular stencils can produce outlying
gradients --- and inactive in smooth regions, so the measured order is
unaffected. The polyhedral studies additionally select the compact
face-neighbour stencil (\texttt{stencil face}) for the gradient, the
mesh-appropriate choice when a polyhedron has 12--16 face neighbours.

\subsection{Phase indicator and narrow band}
\label{sec:indicator}
The volume fraction $\alpha$ is computed from $\psilev$ by the
tetrahedral-decomposition phase indicator of Detrixhe and Aslam
\cite{detrixhe2016}: each cell is decomposed into tetrahedra from its
centroid and face centroids, $\psilev$ is interpolated to the tetrahedron
vertices from a linear least-squares reconstruction, and the wet volume of
each tetrahedron is integrated exactly from the vertex values' sign
pattern. The indicator is evaluated in a narrow band around the interface
($\alpha$ is $0/1$ outside), and all band-restricted error measures of
Section~\ref{sec:metrics} use this band. The indicator and band machinery
are shared verbatim with the quadratic companion line.

% ============================================================================
\section{Software design and reproducibility}
\label{sec:software}

The reconstruction is a run-time-selectable strategy behind the
\texttt{slReconstruction} interface: the case dictionary entry
\begin{center}
\texttt{levelSet \{ semiLagrangian \{ reconstruction nestedLSQ; \} \}}
\end{center}
is the \emph{only} difference between the linear and the quadratic study ---
the foot, the phase indicator, the band, the error function objects and the
whole study pipeline are shared, and swapping \texttt{nestedLSQ} for
\texttt{linearTaylor} (or the quadratic \texttt{uncachedQuadraticWeighted\-Least\-Squares})
is a one-word change. The method line reproduces with one command,
\begin{center}
  \texttt{make sl-linear}
\end{center}
which runs the five convergence studies
(\texttt{config/linearConv*.yaml}: 2D vortex hex; 3D shear and deformation,
hexahedral and polyhedral), aggregates the error tables, regenerates every
figure and table of this article, mirrors them into the companion
presentation's data folder (so slides and paper are always consistent), and
rebuilds both documents. Completing any \emph{single} study auto-refreshes
the documents as well; the meta-workflow only re-harvests.

% ============================================================================
\section{Verification}
\label{sec:verification}

\subsection{Benchmarks and error measures}
\label{sec:metrics}

\paragraph{Benchmarks.} The 2D reversed (single) vortex of Rider and Kothe
\cite{rider1998}: a circle of radius $0.15$ centred at $(0.5,0.75)$ in the
unit box, sheared by the vortex stream function and reversed at $T/2$ so the
exact solution at $t=T=2$ is the initial circle; resolutions
$N = 32 \dots 256$ on a $\sqrt2$ ladder, at $\mathrm{CFL}=0.5$ and $1.0$.
The 3D shear \cite{leveque1996,liovic2006} and 3D deformation
\cite{leveque1996,enright2002} benchmarks in the unit cube with $T=3$,
$\mathrm{CFL}=0.5$, on hexahedral ladders ($N = 32,50,80,128(,160)$, ratio
$\sqrt[3]{4}$) and on polyhedral (cfMesh \cite{juretic_cfmesh}) ladders with
matching characteristic sizes $h = \texttt{maxCellSize} =
1/32 \dots 1/128$; for the polyhedral meshes $h$ is the direct analogue of
$1/N$, so hex and poly lines are plotted on the same axes.

\paragraph{Error measures.} All errors are reported at the end time $t=T$
(reversed state, exact solution known) and the key ones additionally at the
half time $t=T/2$ (maximally deformed state):
\begin{itemize}
\item \emph{shape error}: the geometric $L_1$ distance between the computed
  and exact volume-fraction fields,
  $E_{\mathrm{shape}} = \sum_c |\alpha_c - \alpha_c^{\mathrm{exact}}|\,V_c$
  (at $t=T$ the exact field is the initial one);
\item \emph{volume-conservation error}: the relative signed residual
  $E_{\mathrm{vol}} = (\sum_c \alpha_c V_c - V^0)/V^0$; semi-Lagrangian
  advection is non-conservative by construction, so this is a diagnostic
  bound, not a telescoping conservation identity --- it is bounded and
  non-monotone, and we plot its magnitude without quoting a fitted order;
\item \emph{gradient (signed-distance) defect}:
  $E_{\nabla\psi} = \| |\nabla\psilev| - 1 \|_{L_2}$, globally and
  restricted to the narrow band (superscript ``band'') that the numerics
  actually use;
\item \emph{half-time variants}: $E_{\nabla\psi}(T/2)$,
  $E^{\mathrm{band}}_{\nabla\psi}(T/2)$ and $E_{\mathrm{vol}}(T/2)$ probe
  the maximally stretched interface, where the level set is furthest from a
  signed-distance field;
\item \emph{health metrics}: $\min|\nabla\psilev|$ in the band (at $T$ and
  $T/2$) monitors flattening --- a value approaching zero signals imminent
  loss of the interface, before any error norm reacts.
\end{itemize}
The complete per-resolution data are in
\texttt{data/tables/lsl\_convergence.csv}; least-squares convergence orders
over each ladder are in \texttt{data/tables/lsl\_convergence\_orders.csv}
and Tables~\ref{tab:orders} and \ref{tab:orders-ext}.

\begin{table}[t]
  \centering
  \caption{Measured least-squares convergence orders of the primary error
  measures (shape, volume, gradient defect at $t=T$), per case, mesh type
  and CFL. Volume-conservation is a signed, bounded residual; its fitted
  slope is reported for completeness only.}
  \label{tab:orders}
  \guardedinput{data/tables/convergence_orders.tex}
\end{table}

\begin{table}[t]
  \centering
  \caption{Measured least-squares convergence orders of \emph{all} error
  measures, including the band-restricted gradient defect and the half-time
  ($t=T/2$, maximum deformation) variants.}
  \label{tab:orders-ext}
  \guardedinput{data/tables/convergence_orders_extended.tex}
\end{table}

\subsection{Two-dimensional convergence}
\label{sec:conv2d}
Figure~\ref{fig:conv2d} shows the primary errors of the reversed 2D vortex
at both CFL numbers, with $O(h)$ and $O(h^{2})$ guides;
Figure~\ref{fig:conv2d-diag} completes the picture with the band-restricted
and half-time measures. Three behaviours stand out. First, at
$\mathrm{CFL}=\tfrac12$ the method converges: the shape error at order
$\approx1.1$, the volume error at $\approx1.5$, and the band gradient defect
at $\approx2.1$ (Table~\ref{tab:orders}) --- monotone over the whole ladder.
Second, the absolute errors sit well above the quadratic line's at equal $h$
--- the price of the missing curvature term --- so the linear line buys
accuracy at first-to-second-order rates rather than the quadratic line's
$\approx3$. Third, and unlike the quadratic line, the method is
\emph{conditionally} stable in the Courant number: at $\mathrm{CFL}=1$ it
converges over a bounded envelope (shape order $0.97$ through $N=90$) but
destabilizes on the finest three levels ($N\ge128$), where the same
gradient-steepening feedback that destroys the \texttt{linearTaylor} variant
(Section~\ref{sec:lintaylor}) begins to bite the least-squares slope as well.
The figure fits and draws only the stable envelope of each series and marks
the first destabilized resolution separately (red cross); the quadratic
reconstruction shows no such CFL sensitivity. Production linear studies ---
including all 3D cases below --- run at $\mathrm{CFL}=\tfrac12$. The half-time
gradient defect (maximum stretching) grows toward $O(1)$ on all meshes for
\emph{both} the linear and the quadratic line: neither reconstruction
maintains the signed-distance property through the maximally deformed state.

\begin{figure}[t]
  \centering
  \guardedgraphics{lsl_2Dvortex_convergence.png}
  \caption{Reversed 2D vortex ($T=2$), least-squares linear (\texttt{nestedLSQ})
  semi-Lagrangian method: gradient defect, shape error and volume-conservation
  error at $t=T$ vs $h=1/N$, $\mathrm{CFL}=0.5$ and $1.0$, with measured
  least-squares orders in the legends. Only the stable envelope of each series
  is drawn and fitted; at $\mathrm{CFL}=1$ the first destabilized resolution
  ($N=128$) is marked separately (red cross), the runaway values kept off the
  axis.}
  \label{fig:conv2d}
\end{figure}

\begin{figure}[t]
  \centering
  \guardedgraphics{lsl_2Dvortex_convergence_diagnostics.png}
  \caption{Reversed 2D vortex, all remaining error measures: band-restricted
  gradient defect at $t=T$, global and band-restricted gradient defect at
  the maximally deformed half time $t=T/2$, and the half-time volume error.}
  \label{fig:conv2d-diag}
\end{figure}

\subsection{Three-dimensional convergence, hexahedral and polyhedral}
\label{sec:conv3d}
Figures~\ref{fig:conv3dshear} and~\ref{fig:conv3ddef} overlay the hexahedral
and polyhedral ladders of the 3D shear and deformation benchmarks (hex solid,
poly dashed); polyhedral runs use the face-neighbour gradient stencil and the
stencil-bounds clamp (Section~\ref{sec:safeguards}). The two benchmarks
separate cleanly, and the separation is the central 3D result.

The 3D \emph{deformation} case (Figure~\ref{fig:conv3ddef}) converges on both
mesh types: shape order $2.28$ (hex, five levels to $N=160$) and $1.52$
(poly), volume $3.6$/$2.9$, and the band gradient defect \emph{decreases} with
refinement --- the same well-behaved picture as the 2D vortex, at higher
order. The 3D \emph{shear} case (Figure~\ref{fig:conv3dshear}) converges on the
hexahedral mesh over a bounded envelope and then destabilizes: it refines
cleanly through $N=50$ (shape order $1.18$ over the stable levels) but by the
finer resolutions the gradient defect has run away (to $\mathcal{O}(10^{5})$ at
$N=128$, where the shape error itself jumps to $3.7\times10^{-2}$). The figure
and Table~\ref{tab:orders} therefore fit and report the hexahedral shear order
over the \emph{stable envelope} only ($h\ge2\times10^{-2}$) and mark the first
destabilized resolution separately as the stability limit; the runaway values
are deliberately kept off the axis. This is the same conditional-stability
limit seen at $\mathrm{CFL}=1$ in 2D (Section~\ref{sec:conv2d}), here reached
at $\mathrm{CFL}=\tfrac12$ because the LeVeque--Liovic shear
\cite{leveque1996,liovic2006} sustains a stronger, non-reversing strain than
the deformation field over most of the period, and that strain steepens
$\nabla\psilev$ until the least-squares slope can no longer damp it. The limit
is therefore \emph{flow-dependent}: the 2D vortex (to $N=256$) and the 3D
deformation (to $N=160$) stay inside it at $\mathrm{CFL}=\tfrac12$; the 3D
shear on hexahedra does not, past $N=50$. Notably the \emph{polyhedral} shear
run converges over its full ladder (shape order $1.78$, band gradient $2.10$):
the face-neighbour stencil's stronger averaging and the stencil-bounds clamp
push the limit past the reachable resolution. The practical rule is that
\texttt{nestedLSQ} is a conditionally stable method whose safe resolution
depends on the flow; for the strongest sustained shear on hexahedra it is
bounded, and either a lower Courant number, the face-neighbour stencil, or
band-limited redistancing is required to go finer.

\begin{figure}[t]
  \centering
  \guardedgraphics{lsl_3Dshear_convergence.png}
  \guardedgraphics{lsl_3Dshear_convergence_diagnostics.png}
  \caption{3D shear ($[0,1]^3$, $T=3$, $\mathrm{CFL}=0.5$), hexahedral and
  polyhedral meshes: primary errors (top) and remaining error measures
  (bottom) vs $h$. Only the stable envelope is drawn and fitted: the
  hexahedral run converges over $h\ge2\times10^{-2}$ (shape order $1.18$) and
  then destabilizes (gradient defect $\mathcal{O}(10^{5})$ by $N=128$), the
  first destabilized resolution marked separately (red cross) with the runaway
  values off the axis; the polyhedral run converges throughout.}
  \label{fig:conv3dshear}
\end{figure}

\begin{figure}[t]
  \centering
  \guardedgraphics{lsl_3Ddeformation_convergence.png}
  \guardedgraphics{lsl_3Ddeformation_convergence_diagnostics.png}
  \caption{3D deformation (unit cube, $T=3$, $\mathrm{CFL}=0.5$), hexahedral
  ($N=32\!\to\!160$) and polyhedral meshes: primary errors (top) and remaining
  error measures (bottom) vs $h$. Both mesh types converge (shape order $2.28$
  hex, $1.52$ poly); the band gradient defect decreases with refinement.}
  \label{fig:conv3ddef}
\end{figure}

\subsection{Volume-fraction field evolution}
\label{sec:fields}
Figure~\ref{fig:alpha2d} shows the volume-fraction field of the 2D vortex at
$t=0$, $T/2$ and $T$ across the resolution ladder ($\mathrm{CFL}=\tfrac12$).
This is the method's behaviour made visible: at $T/2$ the filament tail thins
below the local cell size on the coarse meshes and the graded $\alpha$
transition widens, while the reversed state at $t=T$ shows how much of the
initial circle is recovered (exact circle dashed). The under-resolved tail is
lost irreversibly, exactly as in every interface-capturing method at these
resolutions \cite{enright2002}; the convergence of Section~\ref{sec:conv2d}
quantifies the recovery rate, and the sharpening of the $\alpha$ transition
with refinement is its field-level counterpart.

\begin{figure}[t]
  \centering
  \guardedgraphics{alpha_field_2Dvortex_hex.png}
  \caption{Volume fraction $\alpha$ of the reversed 2D vortex
  ($\mathrm{CFL}=0.5$) at $t=0$, $T/2$, $T$ (rows) across the resolution
  ladder (columns; coarsest, median, finest). Black line: the $\psilev=0$
  level set; dashed red circle on the $t=T$ row: the exact reversed
  interface.}
  \label{fig:alpha2d}
\end{figure}

\subsection{Interface evolution in 3D}
\label{sec:iso3d}
Figures~\ref{fig:iso3dshear} and~\ref{fig:iso3ddef} render the
$\psilev = 0$ iso-surface of the 3D benchmarks at $t=0$, $T/2$ and $T$ on
the three finest resolutions of each ladder (hexahedral and polyhedral), all
at $\mathrm{CFL}=\tfrac12$. The deformation interface recovers cleanly at
$t=T$ on both meshes, and the qualitative signature of the linear method is
the earlier loss of the thin scroll edge than the quadratic companion at equal
resolution --- the first-order transport error made visible. The hexahedral
\emph{shear} montage additionally shows the destabilization of
Section~\ref{sec:conv3d}: its finest column ($N=128$) carries the corrupted
interface of the runaway gradient, whereas the polyhedral shear montage
recovers the sheet at every resolution.

\begin{figure}[t]
  \centering
  \guardedgraphics[width=0.49\linewidth]{isosurface_3Dshear_hex.png}
  \guardedgraphics[width=0.49\linewidth]{isosurface_3Dshear_poly.png}
  \caption{3D shear: $\psilev=0$ iso-surfaces at $t=0$, $T/2$, $T$ (rows) on
  the three finest resolutions (columns); hexahedral (left) and polyhedral
  (right) meshes.}
  \label{fig:iso3dshear}
\end{figure}

\begin{figure}[t]
  \centering
  \guardedgraphics[width=0.49\linewidth]{isosurface_3Ddeformation_hex.png}
  \guardedgraphics[width=0.49\linewidth]{isosurface_3Ddeformation_poly.png}
  \caption{3D deformation: $\psilev=0$ iso-surfaces at $t=0$, $T/2$, $T$
  (rows) on the three finest resolutions (columns); hexahedral (left) and
  polyhedral (right) meshes.}
  \label{fig:iso3ddef}
\end{figure}

\subsection{The three reconstructions: two design points and one cautionary tale}
\label{sec:comparison}
The framework's reconstructions occupy three positions, not two. The quadratic
UQWLSR is the kinematic-accuracy champion: it carries the curvature term the
linear fits drop, its 2D shape order is $\approx3$ against the linear
$\approx1.1$, and its order is CFL-independent up to the stability bound.
The least-squares linear \texttt{nestedLSQ} is the low-cost, memory-light
transport method documented here: convergent at $\mathrm{CFL}\le\tfrac12$,
$O(1)$ memory per cell, no fine-mesh ceiling --- the choice when accuracy is
adequate at first-to-second order and the mesh is large. The raw Taylor
\texttt{linearTaylor} is the cautionary tale: unstable in pure advection
(Section~\ref{sec:lintaylor}), it must never be used for
reinitialization-free transport.

The coupling to two-phase flow reorders these, and it is worth being precise
about how, because it is the origin of the linear line. There, cell-centred
curvature estimated from the level set feeds surface-tension forces back into
the flow; a reinitialization-free, high-order transport provides \emph{zero}
damping for the level-set perturbations this feedback loop injects, and the
loop's gain grows with resolution. In that setting the dissipative
\texttt{linearTaylor} transport --- unstable on its own over a long reversible
integration --- is measured to \emph{outlive} the quadratic pipeline at every
resolution and to break its finer-blows-first blow-up trend, because the
coupled solver's own dissipation and the short pre-equilibrium time horizon
keep it inside its stable regime while its damping suppresses the feedback the
quadratic transport amplifies. The two facts coexist without contradiction
(Section~\ref{sec:lintaylor}). The practical recipe of the framework is
therefore: quadratic transport for kinematic studies and
interface-accuracy-critical work; \texttt{nestedLSQ} for low-cost stable
kinematic transport at $\mathrm{CFL}\le\tfrac12$; and \texttt{linearTaylor}
reserved for the explicit-capillary two-phase solver, where its instability is
held in check and its dissipation is an asset.

% ============================================================================
\section{Limitations and outlook}
\label{sec:limitations}
The method inherits the two structural limitations of the framework. First,
semi-Lagrangian point updates are non-conservative; the volume error is
bounded and converges with resolution but never telescopes to zero, and
applications that require discrete conservation need a conservative
correction on top. Second, neither reconstruction maintains the
signed-distance property through strong stretching (the half-time gradient
defect is $O(1)$ on all meshes); downstream consumers that assume
$|\nabla\psilev|=1$ --- curvature estimators above all --- need either
band-limited redistancing or estimators that tolerate the defect. Third, and
specific to the linear line, the stability is only conditional and
flow-dependent: \texttt{nestedLSQ} converges to $256^2$ (2D vortex) and $N=160$
(3D deformation) at $\mathrm{CFL}=\tfrac12$, but the strongest sustained 3D
shear destabilizes it on hexahedra past $N=80$ --- the same gradient-steepening
feedback that destroys \texttt{linearTaylor}, here merely deferred by the
stencil averaging rather than removed. The quadratic line is CFL- and
flow-independent up to the stability bound. The natural outlook is the
combination measured in the two-phase companion work: dissipative linear
transport for damping the curvature-feedback loop, plus targeted
signed-distance maintenance in the band, which would also lift the shear
resolution limit by preventing the gradient runaway at its source.

% ============================================================================
\section{Conclusions}
\label{sec:conclusions}
The linear semi-Lagrangian level-set method is the low-cost sibling of the
quadratic reconstruction, and getting it right turns on a single choice. The
least-squares linear fit \texttt{nestedLSQ} --- a compact per-cell solve, no
cached pseudo-inverse, $O(1)$ memory, no fine-mesh ceiling --- is a genuine,
convergent transport method: at $\mathrm{CFL}\le\tfrac12$ it converges
monotonically in shape, volume and the band gradient defect on the standard
2D and 3D benchmarks, hexahedral and polyhedral, at first-to-second order. The
naive alternative, the raw single-cell Taylor extrapolation
\texttt{linearTaylor}, is unstable in pure advection --- its gradient defect
grows geometrically under refinement --- and must not be used for
reinitialization-free transport, even though it is precisely the
reconstruction that serves the two-phase solver, where coupled dissipation and
short integration keep it in check. \texttt{nestedLSQ}'s stability is
conditional and flow-dependent (it converges to $256^2$ and $N=160$ but the
strongest sustained 3D shear bounds it on hexahedra at $N=80$), which is the
same steepening feedback deferred rather than removed. Documented as an
independent method line with its own one-command reproduction, the linear
line is both a practical, memory-light kinematic transport method within its
stable envelope and a clean demonstration that stencil averaging, not the
polynomial degree alone, is what makes a semi-Lagrangian reconstruction
stable.

\section*{Code and data availability}
The method, the studies and this document reproduce from the public
repository with \texttt{make sl-linear} (OpenFOAM v2506, snakemake); every
figure and table in this article is pipeline-generated and mirrored
verbatim into the companion presentation.

\bibliographystyle{elsarticle-num}
\bibliography{refs}
\end{document}
